\section{Conclusion}
This thesis developed a threat-aware, two-stage waste detection and classification system that operates on real-world-condition images. A single-class YOLO11s detector localizes the waste item, and a shared-backbone network with four task heads jointly predicts material (ten classes), condition, disposal route, and an ordinal threat level; the heads are fused by a regulation-grounded decision layer with a threat-based safety backstop, and per-head calibration is evaluated. Among ten classifier configurations spanning convolutional and transformer backbones and an ensemble, ConvNeXt-S (ImageNet-22k) was the strongest individual backbone on the primary material task, and the system was evaluated by cross-validation over 457 items with bootstrap confidence intervals.

The cross-validated results are strong and honest: material accuracy 0.832 (95\% confidence interval 0.797--0.864) and macro-F1 0.803, with material calibration error 0.036, threat QWK 0.771, and condition AUROC 0.821; the held-out test model reaches material accuracy 0.840 and macro-F1 0.840 on 50 items. The detector reaches mAP@50 0.992, and end-to-end detection-plus-classification shows no measurable material-accuracy loss versus an oracle crop on those 50 items, at a detection rate of 1.0, although condition AUROC and route macro-F1 decrease slightly. Above all, the routing policy sends 148 of 169 hazardous items to Special (recall 0.876). Including two referrals for human review raises safety recall to 0.888 (95\% confidence interval 0.838--0.933) at coverage 0.961, showing both the value of safety-aware routing and the remaining missed hazards. The principal contribution is modelling disposal threat as a rank-consistent ordinal output coupled to a regulation-grounded decision, with calibration, external validation, and rare-class analysis providing a broader assessment of the system.

\section{Limitations and Threats to Validity}
\label{sec:threats}
The work has clear, openly stated limitations. The dataset comprises 485 items, and composite and hazardous\_chemical remain the hardest material classes. The cross-validated analysis includes 21 items from each class, with F1 scores of 0.513 and 0.634 respectively; limited support and visual ambiguity make reliable recognition difficult. The threat backstop can intercept some material errors, but it does not contain every hazardous miss. The route head is the weakest head (cross-validated macro-F1 0.675), dragged down by the very small Reuse class and by Recycle/Landfill confusion. Finally, the system is appearance-only, limiting its ability to distinguish visually similar classes. Threat level 0 has no evaluation examples, so performance at that level remains unvalidated.

Three threats to validity are acknowledged explicitly. Internal validity: the held-out test split contains 50 items, so its point estimates carry wide implicit intervals; this is mitigated by reporting cross-validation over 457 items with bootstrap confidence intervals and by paired significance tests for the backbone choice. External validity: all items were photographed in one geographic and equipment context, and the TACO transfer result (mAP@50 0.096) quantifies how far the detector generalizes beyond single-item photographs; conclusions therefore apply to single-item capture settings, not to cluttered street scenes. Construct validity: the ordinal threat labels follow the published Litter Impact Index rubric rather than ad-hoc judgement, but any rubric involves interpretation at class boundaries; the routing matrix pins the label definitions to cited regulations to keep this subjectivity bounded and auditable.

\section{Future Work}
Several directions follow naturally from these limitations:
\begin{itemize}
    \item \textbf{Multi-modal disambiguation.} Adding a second signal such as optical character recognition of hazard pictograms, warning text (``danger,'' ``poison''), and resin/recycling codes is a promising way to separate visually identical classes such as chemical versus ordinary containers, and to improve the rare classes.
    \item \textbf{Larger and broader data.} Collecting more items, especially of the scarce hazardous and composite streams and across more environments, could improve performance and the reliability of rare-class evaluation.
    \item \textbf{Stronger route modelling.} The route head could be improved with route-specific features or by explicitly conditioning on the (calibrated) material and condition predictions.
    \item \textbf{Deployment study.} A field trial on an edge device or sorting station, measuring throughput, latency, and the value of the abstain option with a human in the loop, would assess real-world readiness.
    \item \textbf{Temporal and video input.} Using multiple frames or short video of an item could resolve ambiguities that a single still cannot.
\end{itemize}
Together, these would extend a complete, safety-aware pipeline toward broader, real-world deployment.
